\documentclass[12pt]{article}
\usepackage[utf8]{inputenc}
\usepackage[T1]{fontenc}
\usepackage[margin=1in]{geometry}
\usepackage{setspace}
\usepackage{xcolor}
\usepackage[colorlinks=true, urlcolor=blue!60!black, linkcolor=blue!60!black]{hyperref}
\setlength{\parindent}{0em}
\setlength{\parskip}{0.6em}
\renewcommand{\baselinestretch}{1.0}
\begin{document}
\begin{center}
{\LARGE \textit{United States v. Morrison}, 529 U.S. 598 (2000)\par}
\vspace{0.5em}
{\large Case Brief --- Two-Page Edition\par}
\end{center}

\section{Memory Jogger}
VAWA's civil remedy fell: gender-motivated violence is not economic activity Congress may aggregate, and Section 5 cannot reach private actors (\href{https://www.courtlistener.com/opinion/118363/united-states-v-morrison/}{CourtListener}).

\section{Detailed Case Facts}
Virginia Tech student Christy Brzonkala alleged that football players Antonio Morrison and James Crawford raped her; the university's discipline of Morrison was reduced and then set aside. She sued under 42 U.S.C. \S~13981, which gave victims of ``crimes of violence motivated by gender'' a federal damages action, enacted on a large congressional record.

\section{Procedural History}
The district court held \S~13981 unconstitutional; a Fourth Circuit panel reversed, but the en banc court affirmed.

\section{Judicial Votes}
Affirmed 5--4. Rehnquist wrote for O'Connor, Scalia, Kennedy, and Thomas; Thomas concurred; Souter and Breyer dissented, joined by Stevens and Ginsburg.

\section{Holding}
Congress may not regulate noneconomic violent conduct based on aggregate effects, even with extensive findings. Section 5 reaches only state action, and a nationwide remedy against private individuals is not congruent and proportional.

\section{Analysis of Opinions}
\textbf{Majority:} Gender-motivated violence is ``not, in any sense of the phrase, economic activity''; the findings' but-for chain would let Congress regulate all crime and family law. Under the \textit{Civil Rights Cases}, Section 5 reaches only state action, and the remedy burdened no state official. \textbf{Thomas:} the substantial-effects test should go. \textbf{Souter:} Congress's ``mountain of data'' satisfied rational basis; the majority revived pre-1937 formalism. \textbf{Breyer:} the economic/noneconomic line is unworkable.

\section{Examples: Future Applications}
\textbf{Same-Side:} a federal tort for local bias-motivated assaults fails. \textbf{Opposite-Side:} punishing interstate travel to commit domestic violence is valid. \textbf{Fence-Sitter:} protecting an intrastate endangered species on private land.

\section{Critique}
Randy Barnett praised the defense of enumerated powers, while Thomas wanted more. Robert Post and Reva Siegel argued the Court treated gender violence as ``local'' while monopolizing Fourteenth Amendment enforcement.

\section{Key Quotations}
\begin{quote}
``We accordingly reject the argument that Congress may regulate noneconomic, violent criminal conduct based solely on that conduct's aggregate effect on interstate commerce. The Constitution requires a distinction between what is truly national and what is truly local.''
\end{quote}
\noindent --- Chief Justice Rehnquist, Opinion of the Court, 529 U.S. 598.

\end{document}
